% BAWM Paper 1 — Bayesian Assembly as World Modeling
% Proof of concept on synthetic genome
\documentclass[11pt,a4paper]{article}

% ── Packages ─────────────────────────────────────────────
\usepackage[utf8]{inputenc}
\usepackage[T1]{fontenc}
\usepackage{amsmath,amssymb,amsthm}
\usepackage{mathtools}
\usepackage{booktabs}
\usepackage{graphicx}
\usepackage[margin=1in]{geometry}
\usepackage{hyperref}
\usepackage{cleveref}
\usepackage{algorithm}
\usepackage{algpseudocode}
\usepackage{natbib}
\usepackage{xcolor}
\usepackage{caption}
\usepackage{subcaption}

\bibliographystyle{plainnat}

% ── Math shortcuts ───────────────────────────────────────
\newcommand{\bx}{\mathbf{x}}
\newcommand{\bz}{\mathbf{z}}
\newcommand{\bmu}{\boldsymbol{\mu}}
\newcommand{\Real}{\mathbb{R}}
\newcommand{\E}{\mathbb{E}}
\newcommand{\KL}{\mathrm{KL}}
\newcommand{\BF}{\mathrm{BF}}
\DeclareMathOperator{\logit}{logit}
\DeclareMathOperator{\sigmoid}{\sigma}

\newtheorem{definition}{Definition}
\newtheorem{proposition}{Proposition}

% ══════════════════════════════════════════════════════════
\title{Bayesian Assembly as World Modeling:\\
  Calibrated Overlap Inference via Analytic Likelihood Ratios}

\author{Raja Balachander Giddi}
\date{\today}

\begin{document}
\maketitle

% ── Abstract ─────────────────────────────────────────────
\begin{abstract}
De novo genome assembly requires deciding which sequencing reads
originate from overlapping genomic regions --- a problem traditionally
solved by heuristic scoring functions that lack probabilistic
interpretation.  We introduce \textbf{BAWM} (Bayesian Assembly as World
Modeling), a framework that treats overlap detection as Bayesian
posterior inference.  Each candidate overlap receives a posterior
probability computed from analytic likelihood ratios under fitted
genuine and null distributions, producing calibrated confidence
estimates rather than arbitrary scores.

On a 100\,kb synthetic genome with simulated PacBio HiFi reads, BAWM
achieves precision 0.781, recall 0.751, and expected calibration error
(ECE) 0.040 --- demonstrating that the posterior probabilities are
faithful to empirical accuracy.  A contrastive-pretrained read encoder
contributes complementary evidence that improves recall by 10
percentage points without degrading calibration.  We identify the
$k$-mer containment discriminability floor as a structural precision
limit and discuss how overlap-length-conditional likelihoods and larger
genomes can address it in future work.
\end{abstract}

% ══════════════════════════════════════════════════════════
\section{Introduction}
\label{sec:intro}

De novo genome assembly reconstructs a genome sequence from
overlapping sequencing reads without a reference.  The core
computational task is \emph{overlap detection}: determining which
pairs of reads originate from overlapping regions of the genome.
Modern assemblers such as Hifiasm \citep{cheng2021hifiasm}, Canu
\citep{koren2017canu}, and Flye \citep{kolmogorov2019assembly} solve
this with heuristic scoring functions --- typically combining seed
matches, alignment scores, and ad hoc filters.  These scores lack
probabilistic semantics: a score of 150 does not mean anything about
the probability that the overlap is real.

This matters for two reasons.  First, downstream decisions (which
overlaps to trust, when to merge contigs, how to handle repeats)
require knowing \emph{how confident} the assembler is, not just its
ranking.  Second, assembly is inherently sequential --- each new read
updates a growing graph --- and Bayesian inference provides the natural
framework for sequential evidence accumulation.

We propose \textbf{BAWM} (Bayesian Assembly as World Modeling), a
framework that reframes overlap detection as Bayesian posterior
inference.  For each candidate read pair $(i,j)$, BAWM computes:
\begin{equation}
  \label{eq:posterior}
  \logit\, p_{ij}
    = \underbrace{\log \frac{p(\bx_{ij} \mid H_1)}{p(\bx_{ij} \mid H_0)}}_{\text{log Bayes factor}}
    + \underbrace{\logit\, \pi}_{\text{prior}},
\end{equation}
where $H_1$ denotes a genuine overlap, $H_0$ denotes no overlap,
$\bx_{ij}$ are observed features of the pair, and $\pi$ is the prior
probability that a candidate pair is genuine.  The posterior
$p_{ij} = \sigmoid(\logit\, p_{ij})$ is a calibrated probability
estimate.

The key contributions of this paper are:
\begin{enumerate}
  \item A \textbf{two-component Bayes factor model} with fitted
    genuine and null distributions for $k$-mer containment and encoder
    cosine similarity, producing analytic likelihood ratios
    (\Cref{sec:model}).
  \item A \textbf{contrastive-pretrained read encoder} that maps reads
    to a latent space where overlapping reads have higher cosine
    similarity, providing evidence complementary to $k$-mer features
    (\Cref{sec:encoder}).
  \item \textbf{Empirical calibration analysis} showing ECE of 0.040
    on a synthetic genome, demonstrating that the posteriors are
    faithful to empirical accuracy (\Cref{sec:results}).
  \item \textbf{Identification of the containment discriminability
    floor} as a structural precision limit, with clear implications for
    future work (\Cref{sec:discussion}).
\end{enumerate}


% ══════════════════════════════════════════════════════════
\section{Background}
\label{sec:background}

\subsection{Overlap Detection in Assembly}

The overlap--layout--consensus (OLC) paradigm requires identifying all
pairs of reads that share a sufficiently long genomic overlap.
Practical assemblers use seeding strategies --- $k$-mers, minimizers,
or sketch-based indices --- to generate candidate pairs, followed by
alignment or scoring to filter true overlaps from noise.

PacBio HiFi reads (15--25\,kb, $>$Q20 accuracy) enable long overlaps
but introduce the challenge of distinguishing genuine overlaps from
pairs that share sequence similarity due to genomic repeats.

\subsection{Bayesian Overlap Scoring}

Given observed features $\bx_{ij}$ for a candidate pair, the posterior
probability of a genuine overlap is:
\begin{equation}
  p(H_1 \mid \bx_{ij})
    = \frac{p(\bx_{ij} \mid H_1)\, \pi}
           {p(\bx_{ij} \mid H_1)\, \pi + p(\bx_{ij} \mid H_0)\,(1-\pi)}.
\end{equation}
On the logit scale, this decomposes additively into log Bayes factor
plus log prior odds (\Cref{eq:posterior}).  This decomposition is
attractive because independent features contribute additively ---
each new source of evidence simply adds its own log Bayes factor to
the running logit.


% ══════════════════════════════════════════════════════════
\section{Model}
\label{sec:model}

BAWM processes reads sequentially, maintaining a probabilistic
assembly graph $\mathcal{G} = (V, E)$ where nodes represent reads and
edges carry posterior overlap probabilities.  For each incoming read,
the system: (1) generates candidate overlaps via minimizer filtering,
(2) ranks candidates by KL-symmetric distance in latent space, (3)
computes per-edge posterior probabilities via analytic Bayes factors,
and (4) updates the graph.

\subsection{Candidate Generation}
\label{sec:candidates}

Each read sequence is converted to canonical $k$-mer tokens
($k=11$, stride 5) via polynomial hashing into a vocabulary of size
$V = 16{,}384$.  Minimizers (window $w=5$) are extracted and stored in
an inverted index.  For a new read, all nodes sharing $\geq 2$
minimizers form the initial candidate set.

\subsection{Latent Space Ranking}

A stochastic read encoder (\Cref{sec:encoder}) maps each read to a
variational embedding $(\bmu_i, \log\boldsymbol{\sigma}^2_i) \in
\Real^{d} \times \Real^{d}$ with $d = 64$.  Candidates are ranked by
symmetrized KL divergence between their variational posteriors:
\begin{equation}
  \label{eq:kl-sym}
  D_{\KL}^{\mathrm{sym}}(i, j)
    = \frac{1}{2}\bigl[D_{\KL}(q_i \| q_j) + D_{\KL}(q_j \| q_i)\bigr],
\end{equation}
where $q_i = \mathcal{N}(\bmu_i, \mathrm{diag}(\boldsymbol{\sigma}^2_i))$.
The top $K = 50$ candidates by smallest $D_{\KL}^{\mathrm{sym}}$
proceed to Bayes factor scoring.

\subsection{Two-Component Bayes Factor}
\label{sec:bayes-factor}

For each candidate pair, we compute two features: $k$-mer containment
$J_{ij}$ and encoder cosine similarity $c_{ij}$.

\paragraph{$k$-mer containment.}
Given the minimizer sets $M_i$ and $M_j$:
\begin{equation}
  J_{ij} = \frac{|M_i \cap M_j|}{\min(|M_i|, |M_j|)}.
\end{equation}
This measures the fraction of shared $k$-mers relative to the smaller
set, reflecting the expected containment under overlap.

\paragraph{Encoder cosine similarity.}
\begin{equation}
  c_{ij} = \frac{\bmu_i \cdot \bmu_j}{\|\bmu_i\| \, \|\bmu_j\|}.
\end{equation}

\paragraph{Likelihood model.}
Both features are modeled under separate genuine ($H_1$) and null
($H_0$) distributions:
\begin{align}
  J_{ij} \mid H_1 &\sim \mathrm{Beta}(\alpha_1, \beta_1), &
  J_{ij} \mid H_0 &\sim \mathrm{Beta}(\alpha_0, \beta_0), \\
  c_{ij} \mid H_1 &\sim \mathcal{N}(\mu_1^c, \sigma_1^c), &
  c_{ij} \mid H_0 &\sim \mathcal{N}(\mu_0^c, \sigma_0^c).
\end{align}
The log Bayes factor combines both features:
\begin{equation}
  \label{eq:log-bf}
  \log \BF_{ij}
    = \underbrace{\log \frac{p(J_{ij} \mid H_1)}{p(J_{ij} \mid H_0)}}_{\text{containment BF}}
    + w_c \cdot \underbrace{\log \frac{p(c_{ij} \mid H_1)}{p(c_{ij} \mid H_0)}}_{\text{cosine BF}},
\end{equation}
where $w_c \in [0,1]$ is a reliability weight that downweights the
cosine Bayes factor to account for residual correlation with the
containment feature.  The final posterior logit is:
\begin{equation}
  \logit\, p_{ij} = \log \BF_{ij} + \logit\, \pi.
\end{equation}

\paragraph{Calibration pass.}
The distribution parameters $(\alpha_1, \beta_1, \alpha_0, \beta_0,
\mu_1^c, \sigma_1^c, \mu_0^c, \sigma_0^c)$ and prior $\pi$ are
estimated from a single no-gradient pass through all reads using the
trained encoder.  Beta parameters are fit via method of moments;
Gaussian parameters are fit from sample mean and standard deviation.
The prior $\pi$ is estimated as the fraction of genuine pairs among
all top-$K$ candidates encountered during calibration.

\begin{definition}[Calibrated posterior]
  A posterior $p_{ij}$ is \emph{calibrated} if, among all pairs
  assigned posterior $p_{ij} \approx p$, the fraction that are genuine
  overlaps is approximately $p$.  We measure calibration via expected
  calibration error (ECE) with 10 equal-width bins.
\end{definition}

\subsection{Edge Loss}
\label{sec:loss}

Given ground-truth labels $y_{ij} \in \{0, 1\}$, we train the encoder
to minimize a focal binary cross-entropy loss:
\begin{equation}
  \mathcal{L}_{\mathrm{edge}}
    = -\frac{1}{|E|} \sum_{(i,j) \in E}
      w_{ij} \bigl[
        y_{ij} \log p_{ij} + (1 - y_{ij}) \log(1 - p_{ij})
      \bigr],
\end{equation}
where $w_{ij} = (1 - p_{ij})^2$ for genuine pairs and $w_{ij} =
p_{ij}^2$ for null pairs (focal weighting).  The total loss adds a
KL regularizer:
\begin{equation}
  \mathcal{L} = \lambda_1 \mathcal{L}_{\mathrm{edge}}
              + \lambda_2 \mathcal{L}_{\KL},
\end{equation}
with $\mathcal{L}_{\KL}$ using free-bits thresholding
\citep{kingma2016improved} at $\delta = 0.5$ nats per dimension.


% ══════════════════════════════════════════════════════════
\section{Read Encoder}
\label{sec:encoder}

\subsection{Architecture}

The read encoder is a lightweight Transformer that maps a sequence of
canonical $k$-mer tokens to a variational embedding.  The architecture
consists of:
\begin{itemize}
  \item An embedding layer mapping $V + 1$ tokens (including padding)
    to $d_{\mathrm{model}} = 64$ dimensions.
  \item A learnable \texttt{[CLS]} token prepended to the input.
  \item A 3-layer Transformer encoder with 4 attention heads,
    feedforward dimension 256, pre-norm architecture, and dropout 0.1.
  \item Linear projection heads $\bmu = W_\mu h_{\texttt{CLS}}$ and
    $\log\boldsymbol{\sigma}^2 = \mathrm{clamp}(W_\sigma
    h_{\texttt{CLS}}, -8, 2)$.
\end{itemize}
Total parameters: $\sim$1.2M.  The encoder runs on CPU for the
synthetic genome experiments.

\subsection{Contrastive Pretraining}
\label{sec:contrastive}

The encoder is pretrained with supervised InfoNCE
\citep{oord2018representation} using ground-truth overlap labels.  For
each anchor read $a$ with a randomly sampled positive $p$ (a read
with genuine overlap) and $N^{-} = 8$ negatives, the loss is:
\begin{equation}
  \mathcal{L}_{\mathrm{NCE}}
    = -\log \frac{\exp(s_{ap} / \tau)}
                 {\exp(s_{ap} / \tau) + \sum_{n} \exp(s_{an} / \tau)},
\end{equation}
where $s_{ij} = \bmu_i \cdot \bmu_j / \sqrt{d}$ is the scaled
dot-product similarity and $\tau$ is a temperature parameter.

\paragraph{Preventing collapse.}
Contrastive pretraining of variational encoders is prone to
representation collapse, where all embeddings converge to a narrow
region.  We employ several countermeasures:
\begin{enumerate}
  \item \textbf{Spread regularizer:} $\mathcal{L}_{\mathrm{spread}}
    = \mathrm{mean}(|\cos(\bmu_i, \bmu_j)|)$ over all pairs in the
    batch, penalizing uniform alignment.
  \item \textbf{High-temperature warmup:} $\tau = 0.5$ for the first
    10 epochs, annealing to $\tau = 0.1$ over the remaining epochs.
  \item \textbf{Unnormalized $\bmu$:} Using dot products rather than
    cosine similarity in the loss allows the encoder to use embedding
    magnitude as a confidence signal.
  \item \textbf{Orthogonality kick:} The $\mu$-head weights are
    perturbed by $\mathcal{N}(0, 0.01)$ at initialization to break
    symmetry.
  \item \textbf{Minimal KL weight:} $\lambda_{\KL} = 0.01$ during
    pretraining to preserve stochastic structure without constraining
    the mean embeddings.
\end{enumerate}
The total pretraining loss is:
\begin{equation}
  \mathcal{L}_{\mathrm{pretrain}}
    = \mathcal{L}_{\mathrm{NCE}}
    + 0.5 \cdot \mathcal{L}_{\mathrm{spread}}
    + 0.01 \cdot \mathcal{L}_{\KL}.
\end{equation}


% ══════════════════════════════════════════════════════════
\section{Experimental Setup}
\label{sec:experiments}

\subsection{Synthetic Genome}

We generate a 100\,kb circular prokaryotic chromosome with 50\% GC
content and three identical 500\,bp repeat insertions spaced at 25\,kb
intervals.  The genome is generated with a fixed random seed for
reproducibility.

\subsection{Read Simulation}

PacBio HiFi reads are simulated at 30$\times$ coverage, producing
$\sim$200 reads with log-normal length distribution (mean 15\,kb,
SD 3\,kb).  The error model applies substitutions at rate $\varepsilon_s
= 0.003$ and indels at rate $\varepsilon_i = 0.001$, with elevated
indel rates at homopolymer runs $\geq 3$\,bp.  Reads are randomly
assigned to forward or reverse strand.

\subsection{Ground Truth Labels}

A read pair $(i,j)$ is labeled genuine ($y_{ij} = 1$) if the true
genomic overlap exceeds 5\,000\,bp based on the known read positions.
This threshold is set above the $k$-mer containment discriminability
floor, where genuine and null containment distributions overlap
significantly (see \Cref{sec:disc-floor}).

\subsection{Training Protocol}

\begin{enumerate}
  \item \textbf{Contrastive pretraining} (30 epochs): supervised
    InfoNCE with batch size 4, learning rate $3 \times 10^{-4}$,
    cosine annealing schedule.
  \item \textbf{Calibration}: single no-gradient pass to fit the
    two-component Bayes factor model parameters.
  \item \textbf{Evaluation}: frozen encoder with fitted Bayes factor
    model, computing posterior probabilities for all candidate edges.
\end{enumerate}

\subsection{Metrics}

We report precision and recall at the default decision threshold
($p > 0.5$, equivalently $\logit\, p > 0$), and expected calibration
error (ECE) with 10 equal-width bins:
\begin{equation}
  \mathrm{ECE} = \sum_{b=1}^{10} \frac{n_b}{N}
    \left| \mathrm{acc}(b) - \mathrm{conf}(b) \right|,
\end{equation}
where $\mathrm{acc}(b)$ and $\mathrm{conf}(b)$ are the mean accuracy
and mean confidence in bin $b$.


% ══════════════════════════════════════════════════════════
\section{Results}
\label{sec:results}

\subsection{Main Results}

\Cref{tab:main} presents the ablation results.  The containment-only
Bayes factor achieves precision 0.777 and recall 0.650 with ECE 0.045.
Adding the contrastive encoder's cosine similarity as a
reliability-weighted second Bayes factor feature improves recall by
10.1 percentage points to 0.751 while maintaining calibration
(ECE 0.040).

\begin{table}[ht]
  \centering
  \caption{Ablation results on the 100\,kb synthetic genome.
    Bold indicates best result.}
  \label{tab:main}
  \begin{tabular}{lccc}
    \toprule
    Configuration & Precision & Recall & ECE \\
    \midrule
    Containment BF only        & 0.777 & 0.650 & 0.045 \\
    + Contrastive encoder      & \textbf{0.781} & \textbf{0.751} & \textbf{0.040} \\
    \midrule
    Target                     & $>$0.80 & $>$0.75 & $<$0.12 \\
    \bottomrule
  \end{tabular}
\end{table}

\subsection{Calibration Analysis}

The reliability diagram (\Cref{fig:calibration}) shows that the
predicted posterior probabilities closely track empirical accuracy
across all confidence bins.  The ECE of 0.040 indicates that the
Bayesian framework produces posteriors that are faithful to the true
overlap probability --- when BAWM assigns $p = 0.7$ to an edge,
approximately 70\% of such edges are genuine overlaps.

\begin{figure}[ht]
  \centering
  \includegraphics[width=0.9\textwidth]{../results/calibration.png}
  \caption{Left: reliability diagram showing predicted confidence
    vs.\ empirical accuracy.  Points near the diagonal indicate good
    calibration (ECE = 0.040).  Right: confidence distribution by
    ground-truth label.  The model assigns high confidence primarily to
    genuine overlaps and low confidence to null pairs.}
  \label{fig:calibration}
\end{figure}

\subsection{Logit Distribution Analysis}

\Cref{tab:logits} shows the logit distributions for genuine and null
edges.  The genuine distribution is centered at $+1.44$ with 76.4\% of
genuine pairs above the decision threshold, while the null distribution
is centered at $-4.06$ with only 10.3\% leaking above threshold.

\begin{table}[ht]
  \centering
  \caption{Logit distribution statistics.}
  \label{tab:logits}
  \begin{tabular}{lcccc}
    \toprule
    Class & Mean & Std & Above 0 & Count \\
    \midrule
    Genuine & $+1.44$ & 2.12 & 76.4\% & 2785 \\
    Null    & $-4.06$ & 4.43 & 10.3\% & 5940 \\
    \bottomrule
  \end{tabular}
\end{table}

\subsection{Fitted Distribution Parameters}

The calibration pass yields the following two-component model
(\Cref{tab:distributions}):

\begin{table}[ht]
  \centering
  \caption{Fitted two-component model parameters from calibration.}
  \label{tab:distributions}
  \begin{tabular}{llcc}
    \toprule
    Feature & Distribution & Genuine ($H_1$) & Null ($H_0$) \\
    \midrule
    Containment $J$ & Beta & $\mathrm{Beta}(12.57, 38.72)$ & $\mathrm{Beta}(9.27, 52.60)$ \\
    & Mean & 0.245 & 0.150 \\
    \midrule
    Cosine $c$ & Gaussian & $\mathcal{N}(0.819, 0.143)$ & $\mathcal{N}(0.530, 0.338)$ \\
    \midrule
    Prior $\pi$ & & \multicolumn{2}{c}{0.330 ($\logit\, \pi = -0.707$)} \\
    Cosine weight $w_c$ & & \multicolumn{2}{c}{0.5} \\
    \bottomrule
  \end{tabular}
\end{table}

The genuine containment distribution (mean 0.245) is well-separated
from the null (mean 0.150), providing the backbone discriminative
signal.  The encoder cosine similarity shows larger separation
(genuine mean 0.819 vs.\ null mean 0.530) but with higher variance,
particularly in the null class ($\sigma = 0.338$), reflecting the
encoder's uncertainty on pairs without genuine overlap.


% ══════════════════════════════════════════════════════════
\section{Discussion}
\label{sec:discussion}

\subsection{Why Bayes Factors, Not Z-Scores}

An earlier version of BAWM used $z$-score normalization to convert
raw features into edge confidences.  This approach is fundamentally
flawed: the $z$-score compares each candidate to the empirical
\emph{mixture} of genuine and null candidates, creating a structural
symmetry where approximately 50\% of candidates fall above any
mean-based threshold regardless of the true discriminability.

The Bayes factor formulation avoids this by comparing each observation
against \emph{separate} genuine and null models.  A containment value
of 0.20 receives different treatment depending on whether it is more
likely under $\mathrm{Beta}(12.57, 38.72)$ or
$\mathrm{Beta}(9.27, 52.60)$ --- the likelihood ratio captures the
evidential meaning of the observation rather than its position in the
mixture.

\subsection{The Containment Discriminability Floor}
\label{sec:disc-floor}

Precision is bounded at 0.781, falling short of the 0.80 target.  This
is not a model failure but a structural limit of the $k$-mer
containment feature.

For two reads with genomic overlap $L$ bp, the expected containment
depends on $L$ relative to read length.  For overlaps below
approximately 8\,kb (on reads of mean length 15\,kb), the containment
values for genuine pairs overlap significantly with the null
containment distribution.  The threshold sweep (\Cref{tab:threshold})
confirms that no single threshold simultaneously achieves precision
$> 0.80$ and recall $> 0.75$:

\begin{table}[ht]
  \centering
  \caption{Precision--recall at selected decision thresholds.}
  \label{tab:threshold}
  \begin{tabular}{rcc}
    \toprule
    Logit threshold & Precision & Recall \\
    \midrule
    $-0.25$ & 0.769 & 0.793 \\
    $\phantom{-}0.00$ & 0.781 & 0.751 \\
    $+0.50$ & 0.802 & 0.662 \\
    \bottomrule
  \end{tabular}
\end{table}

The posterior is \emph{correctly calibrated} in this ambiguous regime:
pairs near the discriminability floor receive posterior probabilities
around 0.4--0.6, accurately reflecting the model's genuine uncertainty.
This is good Bayesian behavior --- the model knows what it doesn't know.

\subsection{Contribution of the Contrastive Encoder}

The contrastive encoder provides evidence that is partially independent
of $k$-mer containment.  Before pretraining, the encoder produced
nearly identical cosine similarities for genuine and null pairs
(separation 0.003, indicating complete collapse).  After contrastive
pretraining with the anti-collapse measures described in
\Cref{sec:contrastive}, the separation increased to 0.553 (genuine
$\mathcal{N}(0.766, 0.217)$ vs.\ null $\mathcal{N}(0.212, 0.389)$).

This additional evidence rescues genuine pairs that containment alone
finds ambiguous --- primarily shorter overlaps where $k$-mer sharing is
limited but sequence-level patterns are still detectable by the
encoder.  The recall improvement from 0.650 to 0.751 comes almost
entirely from this population.

Precision does not improve because the false positives are pairs with
moderate containment \emph{and} moderate cosine similarity --- cases
where both features are genuinely ambiguous.  The cosine Bayes factor
is reliability-weighted ($w_c = 0.5$) to prevent correlation with
containment from inflating the combined evidence.

\subsection{Why Online Bayesian Training Did Not Help}

We attempted end-to-end Bayesian training where the encoder is updated
via BCE loss on the two-feature posterior.  This degraded results
because the Bayes factor model parameters (Beta shapes, Gaussian
means/variances) are fitted once during calibration and held fixed.
As the encoder's cosine distribution shifts during training, the
calibration-epoch Gaussians become stale, and the cosine Bayes factor
becomes increasingly miscalibrated.

Periodic refitting of the Bayes factor model during training would
address this, but on the 100\,kb toy genome the encoder lacks
sufficient diversity to improve beyond the contrastive pretraining
checkpoint.  This motivates the move to larger genomes in future work.


% ══════════════════════════════════════════════════════════
\section{Related Work}
\label{sec:related}

\paragraph{Overlap detection.}
Modern long-read assemblers use seed-and-extend strategies with
$k$-mer or minimizer indices \citep{li2016minimap2,
cheng2021hifiasm}.  BAWM shares the minimizer filtering stage but
replaces alignment scoring with probabilistic inference.

\paragraph{Learned representations for genomics.}
Foundation models for DNA sequences \citep{ji2021dnabert,
dalla2023nucleotide} learn general-purpose embeddings but are not
designed for the overlap detection task.  Our contrastive pretraining
is task-specific, directly optimizing for overlap discriminability.

\paragraph{Bayesian approaches in bioinformatics.}
Bayesian methods have been applied to variant calling
\citep{garrison2012haplotype}, phylogenetics \citep{ronquist2003mrbayes},
and read error correction \citep{greenfield2014bayeshammer}, but we
are not aware of prior work formulating overlap detection itself as
Bayesian posterior inference with calibrated likelihood ratios.

\paragraph{Calibration in machine learning.}
Post-hoc calibration methods such as Platt scaling
\citep{platt1999probabilistic} and temperature scaling
\citep{guo2017calibration} are widely used but operate as
corrections to miscalibrated models.  BAWM achieves calibration by
construction through proper Bayesian inference.


% ══════════════════════════════════════════════════════════
\section{Limitations}
\label{sec:limitations}

\begin{enumerate}
  \item \textbf{Toy genome.} The 100\,kb synthetic genome is far
    smaller than real prokaryotic ($\sim$5\,Mb) or eukaryotic genomes.
    The repeat structure (three identical 500\,bp copies) does not
    capture the complexity of real repeat families.

  \item \textbf{Supervised labels.} The current framework requires
    ground-truth overlap labels for the calibration pass.  In practice,
    these would need to be derived from alignment or estimated via
    self-supervision.

  \item \textbf{Single overlap threshold.} Labeling all pairs above
    5\,000\,bp as genuine creates a heterogeneous positive class.
    Overlap-length-conditional likelihoods would better model the
    relationship between overlap length and feature distributions.

  \item \textbf{Frozen Bayes factor model.} The distribution
    parameters are fitted once and not updated as the encoder evolves.
    This prevents effective end-to-end training.

  \item \textbf{Minimizer set saturation.} The minimizer count
    (stored as a set) saturates for long overlaps, preventing the
    containment feature from tracking overlap length linearly beyond
    $\sim$5\,kb.
\end{enumerate}


% ══════════════════════════════════════════════════════════
\section{Future Work}
\label{sec:future}

Paper~2 will address the identified limitations:
\begin{itemize}
  \item \textbf{Larger genomes} (1--5\,Mb) with realistic repeat
    families, where the encoder has sufficient diversity to learn
    beyond the contrastive pretraining checkpoint.
  \item \textbf{Online Bayes factor refitting:} periodically
    recomputing the genuine and null distribution parameters as the
    encoder evolves during training.
  \item \textbf{Overlap-length-conditional likelihoods:} separate
    Beta distributions for different overlap-length regimes,
    addressing the heterogeneous positive class.
  \item \textbf{Self-supervised calibration:} deriving pseudo-labels
    from alignment confidence to remove the dependence on ground-truth
    overlap labels.
\end{itemize}


% ══════════════════════════════════════════════════════════
\section{Conclusion}
\label{sec:conclusion}

We have presented BAWM, a framework for de novo genome assembly that
treats overlap detection as Bayesian posterior inference.  On a
synthetic genome, BAWM produces well-calibrated posterior probabilities
(ECE 0.040) that faithfully reflect the true probability of overlap.
The two-component Bayes factor model with $k$-mer containment and
contrastive encoder cosine similarity achieves precision 0.781 and
recall 0.751, with a clear structural explanation for the remaining
precision gap.

The central insight is that proper Bayesian inference --- with fitted
genuine and null distributions and analytic likelihood ratios ---
achieves calibration by construction, without requiring post-hoc
correction.  The model's uncertainty estimates are meaningful: when
BAWM is uncertain about an overlap, it is genuinely ambiguous given
the available evidence.

This proof of concept establishes the theoretical and empirical
foundation for scaling Bayesian assembly to real genomes, where the
framework's calibrated uncertainty estimates can inform downstream
assembly decisions in a principled way.


% ══════════════════════════════════════════════════════════
\bibliography{references}

\end{document}
